\IfFileExists{stacks-project.cls}{%
\documentclass{stacks-project}
}{%
\documentclass{amsart}
}

% For dealing with references we use the comment environment
\usepackage{verbatim}
\newenvironment{reference}{\comment}{\endcomment}
%\newenvironment{reference}{}{}
\newenvironment{slogan}{\comment}{\endcomment}
\newenvironment{history}{\comment}{\endcomment}

% For commutative diagrams we use Xy-pic
\usepackage[all]{xy}

% We use 2cell for 2-commutative diagrams.
\xyoption{2cell}
\UseAllTwocells

% We use multicol for the list of chapters between chapters
\usepackage{multicol}

% This is generally recommended for better output
\usepackage{lmodern}
\usepackage[T1]{fontenc}

% For cross-file-references
\usepackage{xr-hyper}

% Package for hypertext links:
\usepackage{hyperref}

% For any local file, say "hello.tex" you want to link to please
% use \externaldocument[hello-]{hello}
\externaldocument[introduction-]{introduction}
\externaldocument[conventions-]{conventions}
\externaldocument[sets-]{sets}
\externaldocument[categories-]{categories}
\externaldocument[topology-]{topology}
\externaldocument[sheaves-]{sheaves}
\externaldocument[sites-]{sites}
\externaldocument[stacks-]{stacks}
\externaldocument[fields-]{fields}
\externaldocument[algebra-]{algebra}
\externaldocument[brauer-]{brauer}
\externaldocument[homology-]{homology}
\externaldocument[derived-]{derived}
\externaldocument[simplicial-]{simplicial}
\externaldocument[more-algebra-]{more-algebra}
\externaldocument[smoothing-]{smoothing}
\externaldocument[modules-]{modules}
\externaldocument[sites-modules-]{sites-modules}
\externaldocument[injectives-]{injectives}
\externaldocument[cohomology-]{cohomology}
\externaldocument[sites-cohomology-]{sites-cohomology}
\externaldocument[dga-]{dga}
\externaldocument[dpa-]{dpa}
\externaldocument[sdga-]{sdga}
\externaldocument[hypercovering-]{hypercovering}
\externaldocument[schemes-]{schemes}
\externaldocument[constructions-]{constructions}
\externaldocument[properties-]{properties}
\externaldocument[morphisms-]{morphisms}
\externaldocument[coherent-]{coherent}
\externaldocument[divisors-]{divisors}
\externaldocument[limits-]{limits}
\externaldocument[varieties-]{varieties}
\externaldocument[topologies-]{topologies}
\externaldocument[descent-]{descent}
\externaldocument[perfect-]{perfect}
\externaldocument[more-morphisms-]{more-morphisms}
\externaldocument[flat-]{flat}
\externaldocument[groupoids-]{groupoids}
\externaldocument[more-groupoids-]{more-groupoids}
\externaldocument[etale-]{etale}
\externaldocument[chow-]{chow}
\externaldocument[intersection-]{intersection}
\externaldocument[pic-]{pic}
\externaldocument[weil-]{weil}
\externaldocument[adequate-]{adequate}
\externaldocument[dualizing-]{dualizing}
\externaldocument[duality-]{duality}
\externaldocument[discriminant-]{discriminant}
\externaldocument[derham-]{derham}
\externaldocument[local-cohomology-]{local-cohomology}
\externaldocument[algebraization-]{algebraization}
\externaldocument[curves-]{curves}
\externaldocument[resolve-]{resolve}
\externaldocument[models-]{models}
\externaldocument[functors-]{functors}
\externaldocument[equiv-]{equiv}
\externaldocument[pione-]{pione}
\externaldocument[etale-cohomology-]{etale-cohomology}
\externaldocument[proetale-]{proetale}
\externaldocument[relative-cycles-]{relative-cycles}
\externaldocument[more-etale-]{more-etale}
\externaldocument[trace-]{trace}
\externaldocument[crystalline-]{crystalline}
\externaldocument[spaces-]{spaces}
\externaldocument[spaces-properties-]{spaces-properties}
\externaldocument[spaces-morphisms-]{spaces-morphisms}
\externaldocument[decent-spaces-]{decent-spaces}
\externaldocument[spaces-cohomology-]{spaces-cohomology}
\externaldocument[spaces-limits-]{spaces-limits}
\externaldocument[spaces-divisors-]{spaces-divisors}
\externaldocument[spaces-over-fields-]{spaces-over-fields}
\externaldocument[spaces-topologies-]{spaces-topologies}
\externaldocument[spaces-descent-]{spaces-descent}
\externaldocument[spaces-perfect-]{spaces-perfect}
\externaldocument[spaces-more-morphisms-]{spaces-more-morphisms}
\externaldocument[spaces-flat-]{spaces-flat}
\externaldocument[spaces-groupoids-]{spaces-groupoids}
\externaldocument[spaces-more-groupoids-]{spaces-more-groupoids}
\externaldocument[bootstrap-]{bootstrap}
\externaldocument[spaces-pushouts-]{spaces-pushouts}
\externaldocument[spaces-chow-]{spaces-chow}
\externaldocument[groupoids-quotients-]{groupoids-quotients}
\externaldocument[spaces-more-cohomology-]{spaces-more-cohomology}
\externaldocument[spaces-simplicial-]{spaces-simplicial}
\externaldocument[spaces-duality-]{spaces-duality}
\externaldocument[formal-spaces-]{formal-spaces}
\externaldocument[restricted-]{restricted}
\externaldocument[spaces-resolve-]{spaces-resolve}
\externaldocument[formal-defos-]{formal-defos}
\externaldocument[defos-]{defos}
\externaldocument[cotangent-]{cotangent}
\externaldocument[examples-defos-]{examples-defos}
\externaldocument[algebraic-]{algebraic}
\externaldocument[examples-stacks-]{examples-stacks}
\externaldocument[stacks-sheaves-]{stacks-sheaves}
\externaldocument[criteria-]{criteria}
\externaldocument[artin-]{artin}
\externaldocument[quot-]{quot}
\externaldocument[stacks-properties-]{stacks-properties}
\externaldocument[stacks-morphisms-]{stacks-morphisms}
\externaldocument[stacks-limits-]{stacks-limits}
\externaldocument[stacks-cohomology-]{stacks-cohomology}
\externaldocument[stacks-perfect-]{stacks-perfect}
\externaldocument[stacks-introduction-]{stacks-introduction}
\externaldocument[stacks-more-morphisms-]{stacks-more-morphisms}
\externaldocument[stacks-geometry-]{stacks-geometry}
\externaldocument[moduli-]{moduli}
\externaldocument[moduli-curves-]{moduli-curves}
\externaldocument[examples-]{examples}
\externaldocument[exercises-]{exercises}
\externaldocument[guide-]{guide}
\externaldocument[desirables-]{desirables}
\externaldocument[coding-]{coding}
\externaldocument[obsolete-]{obsolete}
\externaldocument[fdl-]{fdl}
\externaldocument[index-]{index}

% Theorem environments.
%
\theoremstyle{plain}
\newtheorem{theorem}[subsection]{Theorem}
\newtheorem{proposition}[subsection]{Proposition}
\newtheorem{lemma}[subsection]{Lemma}

\theoremstyle{definition}
\newtheorem{definition}[subsection]{Definition}
\newtheorem{example}[subsection]{Example}
\newtheorem{exercise}[subsection]{Exercise}
\newtheorem{situation}[subsection]{Situation}

\theoremstyle{remark}
\newtheorem{remark}[subsection]{Remark}
\newtheorem{remarks}[subsection]{Remarks}

\numberwithin{equation}{subsection}

% Macros
%
\def\lim{\mathop{\mathrm{lim}}\nolimits}
\def\colim{\mathop{\mathrm{colim}}\nolimits}
\def\Spec{\mathop{\mathrm{Spec}}}
\def\Hom{\mathop{\mathrm{Hom}}\nolimits}
\def\Ext{\mathop{\mathrm{Ext}}\nolimits}
\def\SheafHom{\mathop{\mathcal{H}\!\mathit{om}}\nolimits}
\def\SheafExt{\mathop{\mathcal{E}\!\mathit{xt}}\nolimits}
\def\Sch{\mathit{Sch}}
\def\Mor{\mathop{\mathrm{Mor}}\nolimits}
\def\Ob{\mathop{\mathrm{Ob}}\nolimits}
\def\Sh{\mathop{\mathit{Sh}}\nolimits}
\def\NL{\mathop{N\!L}\nolimits}
\def\CH{\mathop{\mathrm{CH}}\nolimits}
\def\proetale{{pro\text{-}\acute{e}tale}}
\def\etale{{\acute{e}tale}}
\def\QCoh{\mathit{QCoh}}
\def\Ker{\mathop{\mathrm{Ker}}}
\def\Im{\mathop{\mathrm{Im}}}
\def\Coker{\mathop{\mathrm{Coker}}}
\def\Coim{\mathop{\mathrm{Coim}}}

% Boxtimes
%
\DeclareMathSymbol{\boxtimes}{\mathbin}{AMSa}{"02}

%
% Macros for moduli stacks/spaces
%
\def\QCohstack{\mathcal{QC}\!\mathit{oh}}
\def\Cohstack{\mathcal{C}\!\mathit{oh}}
\def\Spacesstack{\mathcal{S}\!\mathit{paces}}
\def\Quotfunctor{\mathrm{Quot}}
\def\Hilbfunctor{\mathrm{Hilb}}
\def\Curvesstack{\mathcal{C}\!\mathit{urves}}
\def\Polarizedstack{\mathcal{P}\!\mathit{olarized}}
\def\Complexesstack{\mathcal{C}\!\mathit{omplexes}}
% \Pic is the operator that assigns to X its picard group, usage \Pic(X)
% \Picardstack_{X/B} denotes the Picard stack of X over B
% \Picardfunctor_{X/B} denotes the Picard functor of X over B
\def\Pic{\mathop{\mathrm{Pic}}\nolimits}
\def\Picardstack{\mathcal{P}\!\mathit{ic}}
\def\Picardfunctor{\mathrm{Pic}}
\def\Deformationcategory{\mathcal{D}\!\mathit{ef}}
\usepackage{cleveref}
\crefname{originallemma}{original lemma}{original lemmas}
\Crefname{originallemma}{Original lemma}{Original lemmas}
\makeatletter
\newlabel{lemma-family-is-one}{{023X}{source}{}{sourcepointer-1}{}}
\newlabel{lemma-family-is-one@cref}{{[originallemma][0][]023X}{[1][0][]source}{}{}{}}
\makeatother
\crefname{originallemma}{original lemma}{original lemmas}
\Crefname{originallemma}{Original lemma}{Original lemmas}
\makeatletter
\newlabel{lemma-fully-faithful}{{0241}{source}{}{sourcepointer-2}{}}
\newlabel{lemma-fully-faithful@cref}{{[originallemma][0][]0241}{[1][0][]source}{}{}{}}
\makeatother
\crefname{originallemma}{original lemma}{original lemmas}
\Crefname{originallemma}{Original lemma}{Original lemmas}
\makeatletter
\newlabel{lemma-pullback-family}{{02VS}{source}{}{sourcepointer-3}{}}
\newlabel{lemma-pullback-family@cref}{{[originallemma][0][]02VS}{[1][0][]source}{}{}{}}
\makeatother
\crefname{originaldefinition}{original definition}{original definitions}
\Crefname{originaldefinition}{Original definition}{Original definitions}
\makeatletter
\newlabel{sites-definition-morphism-coverings}{{00VT}{source}{}{sourcepointer-4}{}}
\newlabel{sites-definition-morphism-coverings@cref}{{[originaldefinition][0][]00VT}{[1][0][]source}{}{}{}}
\makeatother
\crefname{originallemma}{original lemma}{original lemmas}
\Crefname{originallemma}{Original lemma}{Original lemmas}
\makeatletter
\newlabel{topologies-lemma-disjoint-union-is-fpqc-covering}{{040I}{source}{}{sourcepointer-5}{}}
\newlabel{topologies-lemma-disjoint-union-is-fpqc-covering@cref}{{[originallemma][0][]040I}{[1][0][]source}{}{}{}}
\makeatother
\crefname{originallemma}{original lemma}{original lemmas}
\Crefname{originallemma}{Original lemma}{Original lemmas}
\makeatletter
\newlabel{topologies-lemma-recognize-fpqc-covering}{{03L7}{source}{}{sourcepointer-6}{}}
\newlabel{topologies-lemma-recognize-fpqc-covering@cref}{{[originallemma][0][]03L7}{[1][0][]source}{}{}{}}
\makeatother
\title{Pullback of descent data along an fpqc refinement of a family is fully faithful}
\author{The Stacks Project Authors}
\date{stacks-project:descent.tex:local-snapshot-20261012; GNU FDL 1.2 or later}
\begin{document}
\maketitle
\noindent\textit{Editorial note.} All delivered mathematical text is the author's own, copied from the pinned Stacks Project chapter file descent.tex; no statement or proof step was rewritten, reordered or removed.\par\medskip
\noindent\textit{Editorial note.} The proof displays a commutative diagram of four categories of descent data; it lies inside the counted span and is reproduced verbatim.\par\medskip
\noindent\textit{Editorial note.} The counted spans carry 6 navigation stubs; no bibliography citation occurs in the spans.\par\medskip
\noindent\textit{Editorial note.} The preamble span points at the chapter's own preamble.tex shipped with the frozen source, used verbatim as staged.\par\medskip
\noindent\textit{Allowed standard prerequisites (editorial).}\begin{itemize}
\item Morphisms of families of maps with fixed target and descent data relative to a family. Sites, Definition sites-definition-morphism-coverings, invoked at proof lines 8959-8960, together with the same chapter's Lemma lemma-pullback-family invoked at proof lines 8971 and 9007-9008. Fixes the objects the functor acts on and whose commutativity with the coproduct description is used.
\item A disjoint union of fpqc coverings is an fpqc covering. Topologies, Lemma topologies-lemma-disjoint-union-is-fpqc-covering, invoked at proof lines 8984-8986. Turns the hypothesis on each V\_j into an fpqc covering of V\_j by the corresponding sub-coproduct.
\item Recognition criterion for fpqc coverings, part (3). Topologies, Lemma topologies-lemma-recognize-fpqc-covering part (3), invoked at proof lines 8993-8996. Converts the existence of finite affine covers of affine opens of Y into the statement that the single morphism g is an fpqc covering.
\item Descent data for a family agree with descent data for the associated coproduct morphism, and pullback along an fpqc covering of schemes is fully faithful. Same chapter, Lemma lemma-family-is-one at proof lines 9003-9005 and Lemma lemma-fully-faithful at proof lines 9009-9011. Makes the two vertical arrows of the comparison square equivalences and supplies full faithfulness of the top horizontal arrow.
\end{itemize}\medskip
\ifdefined\hypertarget\hypertarget{sourcepointer-1}{}\fi
\noindent\textit{Source reference.} Original lemma 023X: sources/stacks-descent/pinned-original-source.tex, source lines 8467--8467. Pointer outside the counted span, invoked once for context; the load-bearing argument of the counted statement is the author's proof reproduced in full in the proof span.\par\medskip
\ifdefined\hypertarget\hypertarget{sourcepointer-2}{}\fi
\noindent\textit{Source reference.} Original lemma 0241: sources/stacks-algebraic/pinned-original-source.tex, source lines 2746--2746. Pointer outside the counted span, invoked once for context; the load-bearing argument of the counted statement is the author's proof reproduced in full in the proof span.\par\medskip
\ifdefined\hypertarget\hypertarget{sourcepointer-3}{}\fi
\noindent\textit{Source reference.} Original lemma 02VS: sources/stacks-descent/pinned-original-source.tex, source lines 8570--8570. Pointer outside the counted span, invoked once for context; the load-bearing argument of the counted statement is the author's proof reproduced in full in the proof span.\par\medskip
\ifdefined\hypertarget\hypertarget{sourcepointer-4}{}\fi
\noindent\textit{Source reference.} Original definition 00VT: sources/stacks-chow/tags.txt, source lines 1159--1159. Pointer outside the counted span, invoked once for context; the load-bearing argument of the counted statement is the author's proof reproduced in full in the proof span.\par\medskip
\ifdefined\hypertarget\hypertarget{sourcepointer-5}{}\fi
\noindent\textit{Source reference.} Original lemma 040I: sources/stacks-chow/tags.txt, source lines 5257--5257. Pointer outside the counted span, invoked once for context; the load-bearing argument of the counted statement is the author's proof reproduced in full in the proof span.\par\medskip
\ifdefined\hypertarget\hypertarget{sourcepointer-6}{}\fi
\noindent\textit{Source reference.} Original lemma 03L7: sources/stacks-chow/tags.txt, source lines 4703--4703. Pointer outside the counted span, invoked once for context; the load-bearing argument of the counted statement is the author's proof reproduced in full in the proof span.\par\medskip
\begin{lemma}
\label{lemma-fpqc-refinement-coverings-fully-faithful}
Let $S$ be a scheme.
Let $\mathcal{U} = \{U_i \to S\}_{i \in I}$, and
$\mathcal{V} = \{V_j \to S\}_{j \in J}$,
be families of morphisms with target $S$.
Let $\alpha : I \to J$, $\text{id} : S \to S$ and
$g_i : U_i \to V_{\alpha(i)}$ be a morphism of families
of maps with fixed target, see
Sites, Definition \ref{sites-definition-morphism-coverings}.
Assume that for each $j \in J$ the family
$\{g_i : U_i \to V_j\}_{\alpha(i) = j}$ is an fpqc
covering of $V_j$. Then the pullback functor
$$
\text{descent data relative to }
\mathcal{V}
\longrightarrow
\text{descent data relative to }
\mathcal{U}
$$
of Lemma \ref{lemma-pullback-family} is fully faithful.
\end{lemma}

\begin{proof}
Consider the morphism of schemes
$$
g :
X = \coprod\nolimits_{i \in I} U_i
\longrightarrow
Y = \coprod\nolimits_{j \in J} V_j
$$
over $S$ which on the $i$th component maps into the $\alpha(i)$th component
via the morphism $g_{\alpha(i)}$. We claim that $\{g : X \to Y\}$
is an fpqc covering of schemes. Namely, by
Topologies, Lemma \ref{topologies-lemma-disjoint-union-is-fpqc-covering}
for each $j$ the morphism $\{\coprod_{\alpha(i) = j} U_i \to V_j\}$ is an
fpqc covering. Thus for every affine open $V \subset V_j$
(which we may think of as an affine open of $Y$)
we can find finitely many affine opens
$W_1, \ldots, W_n \subset \coprod_{\alpha(i) = j} U_i$
(which we may think of as affine opens of $X$)
such that $V = \bigcup_{i = 1, \ldots, n} g(W_i)$.
This provides enough affine opens of $Y$ which can be covered by finitely
many affine opens of $X$ so that
Topologies, Lemma \ref{topologies-lemma-recognize-fpqc-covering} part (3)
applies, and the claim follows. Let us write $DD(X/S)$,
resp.\ $DD(\mathcal{U})$ for the category of descent data with respect
to $X/S$, resp.\ $\mathcal{U}$, and similarly for $Y/S$ and $\mathcal{V}$.
Consider the diagram
$$
\xymatrix{
DD(Y/S) \ar[r] & DD(X/S) \\
DD(\mathcal{V}) \ar[u]^{\text{Lemma }\ref{lemma-family-is-one}} \ar[r] &
DD(\mathcal{U}) \ar[u]_{\text{Lemma }\ref{lemma-family-is-one}}
}
$$
This diagram is commutative, see the proof of
Lemma \ref{lemma-pullback-family}.
The vertical arrows are equivalences. Hence the lemma follows from
Lemma \ref{lemma-fully-faithful} which shows the top horizontal arrow
of the diagram is fully faithful.
\end{proof}

\end{document}
